% ============ 12. Visualización para observatorios ============
\section{Visualización y comunicación para observatorios}

\subsection{Qué dice la literatura}

El producto público de un observatorio es su capa visual, y existe literatura
madura al respecto. El concepto de \textbf{visualización narrativa}
\citep{segel2010narrative} formaliza cómo contar historias con datos
(estructuras tipo magazine, flujo descendente, paneles tipo "martini glass"),
exactamente el patrón \emph{Pregunta $\rightarrow$ Hallazgo $\rightarrow$ Caveat}
que el proyecto ya usa en su dashboard. Los principios clásicos de diseño
siguen vigentes (datos-ink, eliminación de ruido visual; Tufte, 1983) y los
marcos modernos de dashboards analíticos (Few, 2009) recomiendan jerarquía
visual, contexto y encadenamiento pregunta-respuesta. La literatura de
\emph{visual analytics} enfatiza la \textbf{interacción como medio de
exploración} (filtros, drill-down) y la \textbf{transparencia del dato}
(mostrar la fuente y los supuestos de cada cifra junto a la visualización).
Para observatorios de política, el desafío no es dibujar más, sino
\textbf{auditar qué se muestra}: cada gráfico debe poder trazarse a su dato y a
su método.

\subsection{Relación con este repositorio}

El proyecto entrega un dashboard Streamlit de 7 pestañas (calidad, trayectoria,
territorio, género/diversidad, redes, productividad, datos crudos) con el
patrón Pregunta $\rightarrow$ Hallazgo $\rightarrow$ Caveat, mapas, heatmaps,
Sankeys y grafos interactivos (Pyvis). La auditoría verificó: el gráfico HHI
del dashboard \textbf{nunca se dibuja} por un bug silencioso (ValueError
capturado); la sección de productividad se degrada sin el dataset de
producción; y el manual ejecutable no es reproducible. Frente a la literatura,
faltan: (a) la \textbf{traza de cada visualización a su dato y método} (botón o
nota "fuente y cálculo"), (b) indicadores de \textbf{incertidumbre} en los
gráficos (intervalos de confianza), y (c) una política de auditoría visual de
los gráficos antes de publicarlos.

\subsection{Mejoras recomendadas}

\begin{enumerate}
  \item Documentar cada visualización con su \textbf{fuente, definición del
        cálculo y limitaciones} (encadenando con el marco de indicadores de la
        sección 3 y con el linaje de la sección 4).
  \item Reparar el bug del HHI y añadir \textbf{incertidumbre} (IC) a los
        gráficos de indicadores principales, siguiendo la práctica de
        visualización honesta de la literatura narrativa \citep{segel2010narrative}.
  \item Mantener el patrón \emph{Pregunta $\rightarrow$ Hallazgo $\rightarrow$
        Caveat} como guía de diseño (ya es una fortaleza) y añadir una
        \textbf{vista de auditoría de datos crudos} que permita verificar cada
        hallazgo desde el propio dashboard (tab 7 ya existe; fortalecerla con
        descarga y con la traza metodológica).
\end{enumerate}
